\section*{Capítulo \ref{ch:g2:balance-equilibrium} --- Equilíbrio}

\begin{solution}{exo:g2:balance-equilibrium:1}
A Terra puxa a luminária para baixo; a escrivaninha a empurra para
cima, exatamente com a mesma intensidade. As duas \omterm{def:g2:pushes-and-pulls:force}{forças} empatam ---
\omterm{def:g2:balance-equilibrium:equilibrium}{equilíbrio} --- e por isso a luminária fica onde está.
\end{solution}

\begin{solution}{exo:g2:balance-equilibrium:2}
Quando os dois times puxam com a mesma intensidade em sentidos
opostos: os puxões se anulam. O empate se desfaz no instante em que
um dos times puxa mais forte que o outro.
\end{solution}

\begin{solution}{exo:g2:balance-equilibrium:3}
Ela fica na horizontal --- pesos iguais a distâncias iguais. Quando
uma das gêmeas se arrasta para mais longe, o lado dela desce: a
distância também conta.
\end{solution}

\begin{solution}{exo:g2:balance-equilibrium:4}
Perto do apoio, com Lina lá na ponta do lado dela. A regra: um
passageiro mais pesado perto do meio empata com um passageiro mais
leve lá longe --- o peso e a distância contam os dois.
\end{solution}

\begin{solution}{exo:g2:balance-equilibrium:5}
O \omterm{def:g2:balance-equilibrium:point}{ponto de equilíbrio} da régua fica no meio dela; o do martelo fica
perto da cabeça pesada. Para conferir: apoie cada um em um dedo só e
procure o ponto em que ele fica na horizontal.
\end{solution}

\begin{solution}{exo:g2:balance-equilibrium:6}
Eles se encontram perto da cabeça da vassoura --- a ponta pesada. O
\omterm{def:g2:balance-equilibrium:point}{ponto de equilíbrio} fica do lado pesado de um objeto.
\end{solution}

\begin{solution}{exo:g2:balance-equilibrium:7}
A pirâmide larga: ela pode se inclinar bastante antes que o peso dela
fique para fora da base larga. A torre estreita tomba com a menor
inclinação.
\end{solution}

\begin{solution}{exo:g2:balance-equilibrium:8}
A vara torna mais lento cada desequilíbrio, dando à equilibrista
tempo de voltar para cima da corda. No meio-fio você abre os braços
pelo mesmo motivo: desequilíbrios mais lentos, mais fáceis de
corrigir.
\end{solution}

\begin{solution}{exo:g2:balance-equilibrium:9}
Para ficar em pé só no pé direito, o seu corpo precisa se inclinar
para a direita, de modo que o seu peso fique por cima daquele pé ---
como a torre que mantém o peso por cima da base. A parede está
exatamente onde você precisa se inclinar: ela proíbe o movimento, o
seu peso continua pendurado sobre o lugar vazio do pé levantado, e
você tomba.
\end{solution}

\begin{solution}{exo:g2:balance-equilibrium:10}
O guindaste é uma gangorra gigante, com a torre dele como apoio.
Passageiro um: a carga, lá longe no braço dianteiro comprido.
Passageiro dois: os blocos de concreto --- muito mais pesados, então
podem ficar perto, no braço traseiro curto. Pesado-perto empata com
leve-longe, exatamente a regra da gangorra; o braço de trás é curto
justamente porque o passageiro dele é o peso-pesado.
\end{solution}
